\documentclass[11pt]{article}

\usepackage{amsmath,amssymb,graphicx}

\title{Gyro Logic: A Stability-Based Dynamic Logic Framework}
\author{}
\date{}

\begin{document}

\maketitle

\begin{abstract}
We propose Gyro Logic, a novel logical framework in which meaning,
truth, and inference emerge from the dynamics of behavior fields.
Instead of static symbolic representations, entities are modeled
as trajectories in a behavior manifold.

We introduce a stability functional that governs the persistence
and coherence of behavior under a given frame.
Truth is defined as a stability-weighted projection,
and inference is formalized as a transformation that maximizes stability.

We further define attractors as stability-maximizing regions,
interpreted as meaning cores.
\end{abstract}

\section{Introduction}

Traditional logic assumes static truth values.
However, real-world reasoning is dynamic, contextual, and observer-dependent.

\section{Behavior Representation}

\[
\mathcal{B} : \mathbb{T} \rightarrow \mathcal{M}
\]

\section{Stability Functional}

\[
\mathcal{L}_{stab}[\mathcal{B};F]
=
\alpha C - \beta D - \gamma I + \delta R
\]

\section{Truth}

\[
T(t;F)=\Pi(\mathcal{B},F,t)\cdot\Sigma(\mathcal{B},F,t)
\]

\section{Attractor}

\[
\mathcal{A}_F = \arg\max \mathcal{L}_{stab}
\]

\section{Inference}

\[
\Phi = \arg\max \mathcal{L}_{stab}
\]

\section{Dynamics}

\[
\frac{d\mathcal{B}}{dt} = \nabla \mathcal{L}_{stab}
\]

\section{Conclusion}

Gyro Logic provides a dynamic, stability-based foundation for reasoning.

\end{document}